% =====================================================================
%  Stage 1, раздел 1.
%  Обзор IDS/IPS/NIDS, методов обнаружения аномалий,
%  нейросетевых подходов и метрик качества.
% =====================================================================

\section{Корпоративная сеть как объект мониторинга и обнаружения несанкционированных действий}

\subsection{Понятие корпоративной сети и объекты мониторинга}

В рамках настоящего исследования под корпоративной сетью понимается совокупность
взаимосвязанных программно-аппаратных средств, обеспечивающих обмен данными между
рабочими станциями, серверами, сетевым оборудованием и внешними сегментами (Internet,
филиалы, облачные провайдеры). С точки зрения мониторинга безопасности
корпоративная сеть характеризуется тремя признаками, существенными для построения
системы обнаружения вторжений:

\begin{enumerate}
    \item \textbf{Высокой неоднородностью трафика.} На периметре и в сегментах LAN
    одновременно сосуществуют десятки протоколов прикладного уровня, что приводит к
    мультимодальному распределению признаков сетевых соединений и затрудняет
    применение единых пороговых правил~\cite{khraisat2019survey, ahmad2021nids}.
    \item \textbf{Доминированием легитимного трафика.} Доля атакующих соединений в
    реальной эксплуатации, как правило, не превышает долей процента; это создаёт
    ярко выраженный дисбаланс классов и приводит к необходимости использовать
    специализированные методы обучения и оценки~\cite{sommer2010outside,
    chawla2002smote, lin2017focal}.
    \item \textbf{Динамичностью характеристик.} Состав сетевых сервисов, схемы
    маршрутизации, профили активности пользователей и спектр атак изменяются во
    времени, что приводит к явлению дрейфа распределений (\eng{concept drift,
    covariate shift})~\cite{gama2014drift, andresini2021insomnia}.
\end{enumerate}

В качестве объектов мониторинга в работе рассматриваются: сетевые соединения
(\eng{flows}) на уровне L3/L4, сессии прикладного уровня и временные
последовательности соединений в скользящих окнах. Это согласуется с подходом,
принятым в большинстве современных систем сетевого анализа трафика
(\eng{Network Traffic Analysis, NTA}) и в открытых источниках
данных~\cite{moustafa2015unsw, sharafaldin2018toward}.

\subsection{Модели угроз и таксономия атак}

При проектировании NIDS целесообразно опираться на формализованные таксономии угроз.
Наиболее распространёнными являются следующие модели.

\begin{itemize}
    \item \textbf{MITRE ATT\&CK Enterprise} --- матрица тактик и техник атак,
    систематизирующая действия атакующего на стадиях разведки, первоначального
    доступа, закрепления, эскалации привилегий, латерального перемещения,
    сбора и эксфильтрации данных~\cite{mitreattack}.
    \item \textbf{Таксономия датасета UNSW-NB15}, выделяющая девять семейств
    атакующего трафика: \emph{Fuzzers, Analysis, Backdoors, DoS, Exploits, Generic,
    Reconnaissance, Shellcode, Worms}~\cite{moustafa2015unsw}.
    \item \textbf{Таксономия датасета CICIDS2017}, содержащая категории
    \emph{Brute Force, Heartbleed, Botnet, DoS, DDoS, Web Attack, Infiltration,
    Port Scan}~\cite{sharafaldin2018toward}.
\end{itemize}

В диссертационной работе в качестве рабочей таксономии принята классификация
UNSW-NB15, поскольку именно она используется для разметки основного обучающего
корпуса. Расширение на дополнительные классы атак (\textit{Brute Force, Web
Attack, Botnet}) выполняется на этапе кросс-валидации с CICIDS2017.

\subsection{Виды несанкционированной активности и их сетевые проявления}

С позиции NIDS все рассматриваемые типы вторжений интерпретируются через
проявления в признаках сетевых соединений. Сводная классификация
сетевых проявлений приведена в \tabref{tab:traffic-anomalies}.

\begin{table}[H]
\centering
\caption{Сетевые проявления типовых атак в признаках flow-уровня}
\label{tab:traffic-anomalies}
\renewcommand{\arraystretch}{1.2}
\begin{tabularx}{\linewidth}{@{}lXX@{}}
\toprule
\textbf{Класс атаки} & \textbf{Характеристики во flow-признаках} & \textbf{Признаки временной структуры} \\
\midrule
DoS / DDoS &
Резкий рост числа соединений к одному адресату; короткие потоки; рост доли
неуспешных TCP-handshake; всплески объёма пакетов в единицу времени &
Сжатие межпакетных интервалов, синхронизированные всплески входящего трафика \\
\addlinespace
Reconnaissance, Port Scan &
Большое число потоков к одному источнику с разными destination-портами;
короткие соединения; преобладание SYN/ICMP &
Регулярные интервалы между потоками, монотонный обход диапазона портов \\
\addlinespace
Brute Force &
Повторные подключения к одному порту (22, 3389, 80); ошибки авторизации в
прикладной части; малое число байт за сессию &
Высокая периодичность в окне порядка единиц секунд \\
\addlinespace
Exploits, Backdoors &
Аномальные размеры payload; нестандартные сочетания TCP-флагов;
несвойственное соотношение в обе стороны &
Резкие переходы между фазой сканирования и фазой эксплуатации \\
\addlinespace
Data Exfiltration &
Длинные потоки с высоким \(b_{out}\) и низким \(b_{in}\); устойчивая
скорость передачи; нестандартные протоколы &
Долговременные сессии с малой вариативностью скорости \\
\addlinespace
Worm, Botnet, Lateral movement &
Большое число направленных потоков от одного источника; повторение шаблонных
сессий; обращения к C\&C-инфраструктуре &
Циклическая активность с фиксированным периодом, синхронные всплески \\
\bottomrule
\end{tabularx}
\end{table}

Из \tabref{tab:traffic-anomalies} видно, что значительная часть атак проявляется
не в отдельных признаках, а в их совместном распределении и временной структуре.
Это аргументирует использование моделей, способных одновременно учитывать
локальные паттерны (свёрточные слои) и временные зависимости (рекуррентные слои
или механизмы внимания)~\cite{kim2016lstm, yin2017dlrnn, vinayakumar2019dl}.

\section{Системы обнаружения вторжений: классификация и принципы построения}

\subsection{Классификация IDS по принципу детекции}

В соответствии с устоявшейся в литературе типологией~\cite{khraisat2019survey,
ahmad2021nids, ferrag2020dlcyber} системы обнаружения вторжений делятся на три
крупных класса.

\begin{enumerate}
    \item \textbf{Сигнатурные (\eng{signature-based, misuse detection}).}
    Принимают решение по совпадению наблюдаемого поведения с заранее заданным
    шаблоном (сигнатурой). Преимущества: высокая точность по известным угрозам,
    интерпретируемость, низкая частота ложных срабатываний. Ограничения:
    неприменимость к ранее не встречавшимся атакам (\eng{zero-day}),
    необходимость поддержки баз сигнатур, уязвимость к обфускации трафика.
    Промышленные представители: Snort, Suricata, Zeek (в части сигнатурных
    модулей).
    \item \textbf{Аномальные (\eng{anomaly-based}).} Оценивают отклонение
    наблюдаемого поведения от модели нормы и сигнализируют при превышении
    порога. Преимущества: способность детектировать ранее неизвестные атаки.
    Ограничения: повышенная частота ложных срабатываний, чувствительность к
    дрейфу нормального профиля и к качеству обучающей выборки. Подкласс
    \eng{specification-based} использует формальные модели нормы.
    \item \textbf{Гибридные.} Комбинируют сигнатурную и аномальную ветви,
    распределяя решения по уровню риска и стоимости ошибок: сигнатурная ветвь
    подтверждает известные классы, аномальная --- расширяет покрытие на
    неизвестные сценарии.
\end{enumerate}

В настоящем исследовании рассматривается аномальный класс с акцентом на
нейросетевые модели; гибридная конфигурация выступает рекомендуемой формой
интеграции в SOC, описанной в \secref{subsec:soc-integration}.

\subsection{Packet-based и flow-based анализ}

С точки зрения исходного представления данных различают два уровня анализа.

\textbf{Packet-based (DPI, Deep Packet Inspection).} Объект анализа --- отдельный
сетевой пакет вместе с его полезной нагрузкой. Преимущества: максимальный объём
семантической информации, возможность распознавать атаки на прикладном уровне.
Ограничения: высокая ресурсоёмкость, сложность шифрованного трафика, ограничения
по приватности, затрудняющие использование в корпоративной среде.

\textbf{Flow-based.} Объект анализа --- сетевое соединение (flow), описываемое
агрегированными характеристиками: длительность, количество пакетов и байтов в
обоих направлениях, флаги TCP, межпакетные интервалы, протокол и сервис.
Стандартными протоколами экспорта flow-записей являются Cisco
NetFlow~\cite{rfc3954} и IETF IPFIX~\cite{rfc7011}. Преимущества подхода:
ресурсная эффективность, устойчивость к шифрованию (анализ метаданных), удобство
агрегирования в скользящие окна.

В работе используется flow-based подход. Обоснование: датасеты UNSW-NB15 и
CICIDS2017 представлены именно в этом формате; flow-уровень напрямую
соответствует выходу типовых сборщиков NetFlow/IPFIX, что обеспечивает
переносимость методики в реальные инфраструктуры~\cite{moustafa2015unsw,
sharafaldin2018toward, gharib2016evalframework}.

\subsection{Архитектура SOC и роль NIDS/NTA в общем pipeline алертирования}
\label{subsec:soc-integration}

Современный SOC реализуется как многоуровневый конвейер обработки телеметрии:
сборщики (\eng{network sensors, EDR, log forwarders}) $\rightarrow$ потоковые
очереди (\eng{Kafka, NATS}) $\rightarrow$ нормализаторы $\rightarrow$ хранилища
(SIEM, data lake) $\rightarrow$ аналитические подсистемы (NIDS, UEBA, корреляторы)
$\rightarrow$ модули реагирования (SOAR). Сетевой детектор аномалий выступает
аналитической подсистемой и обязан удовлетворять трём ограничениям:

\begin{itemize}
    \item ограниченное время отклика (\eng{time-to-alert}), достаточное для
    реактивного блокирования;
    \item низкая частота ложных срабатываний на единицу времени
    (\eng{false positives per time}), поскольку каждое срабатывание формирует
    задание для аналитика SOC;
    \item возможность приоритизации алертов, чтобы аналитик в первую очередь
    рассматривал инциденты с высоким влиянием.
\end{itemize}

Указанные ограничения определяют требования к метрикам, рассматриваемым далее в
\secref{sec:metrics}, и к схеме формирования алертов прототипа, описываемой на
этапе 3 настоящей работы.

\section{Методы обнаружения аномалий}

В этом разделе систематизирована теоретическая база основных групп методов,
применяемых для детекции аномалий, и описаны их математические основы.

\subsection{Постановка задачи обнаружения аномалий}

Пусть \(X \subseteq \mathbb{R}^d\) --- пространство признаков сетевых соединений
(или окон), \(p_{0}(x)\) --- плотность распределения нормальных соединений.
Задача состоит в построении функции скоринга \(s\colon X \to \mathbb{R}\) и
порога \(\tau \in \mathbb{R}\) такими, что
\begin{equation}
\hat{y}(x) = \mathbb{I}\bigl[s(x) > \tau\bigr],
\quad x \in X,
\label{eq:detector}
\end{equation}
где \(\hat{y}(x) = 1\) интерпретируется как ``аномалия / атака''.
Обучение проводится в одном из трёх режимов:

\begin{itemize}
    \item \textbf{Supervised}: дана выборка \(\{(x_i, y_i)\}_{i=1}^{n}\) с
    метками классов \(y_i \in \{0,1\}\); скоринг строится через минимизацию
    эмпирического риска.
    \item \textbf{Semi-supervised (one-class)}: даны только нормальные
    наблюдения; скоринг строится как мера отклонения от модели нормы.
    \item \textbf{Unsupervised}: выборка немаркирована; скоринг получается из
    статистических, плотностных или структурных свойств данных.
\end{itemize}

Каждый из режимов применим в условиях NIDS и обладает собственными
ограничениями~\cite{khraisat2019survey, sommer2010outside}. В дальнейшем
исследовании используется преимущественно supervised-режим (на UNSW-NB15) с
параллельным включением semi-supervised-моделей на базе автоэнкодеров и
one-class классификаторов в качестве сравнения.

\subsection{Статистические и энтропийные методы}

Простейший статистический подход использует параметрическое описание
распределения признаков. Например, для одномерного признака \(x\) типичной
является \(z\)-оценка
\begin{equation}
z = \frac{x - \mu}{\sigma},
\qquad
\hat{y} = \mathbb{I}\bigl[\,|z| > \tau_z\bigr],
\label{eq:zscore}
\end{equation}
где \(\mu\), \(\sigma\) --- оценки математического ожидания и
среднеквадратического отклонения. Многомерное обобщение --- расстояние
Махаланобиса
\begin{equation}
d_M(x) = \sqrt{(x-\mu)^{\top}\Sigma^{-1}(x-\mu)},
\label{eq:mahalanobis}
\end{equation}
где \(\Sigma\) --- ковариационная матрица. Подход эффективен при
гауссовости признаков и небольшой размерности, но в реальном NIDS-контексте
сталкивается с многомодальностью распределений и тяжёлыми хвостами.

Энтропийные методы оценивают энтропию распределений ключевых признаков
(распределение IP-адресов, портов, размеров пакетов) в скользящем окне,
интерпретируя резкое изменение энтропии как аномалию. Для дискретного
распределения \(P\) над алфавитом \(\mathcal{A}\) энтропия Шеннона определяется
как
\begin{equation}
H(P) = -\sum_{a \in \mathcal{A}} P(a)\log P(a).
\label{eq:entropy}
\end{equation}
Подобные индикаторы исторически использовались в подсистемах NetFlow-анализа,
но в современных условиях, как отмечается в обзорах~\cite{khraisat2019survey,
ferrag2020dlcyber}, их разрешающая способность недостаточна для распределённых
и низкочастотных атак.

\subsection{Кластеризация и плотностные методы}

\textbf{K-means} формирует разбиение на \(K\) кластеров с центроидами
\(\{c_k\}\) минимизацией внутрикластерной дисперсии:
\begin{equation}
\min_{\{c_k\}}\sum_{k=1}^{K}\sum_{x \in C_k}\|x - c_k\|^{2}.
\label{eq:kmeans}
\end{equation}
Аномалия определяется через расстояние до ближайшего центроида.

\textbf{DBSCAN} определяет аномалии как точки, не принадлежащие ни одному
плотностному кластеру. Параметры: радиус окрестности \(\varepsilon\) и
минимальное число точек \(\mathit{minPts}\). Метод устойчив к
неэллипсоидальной форме кластеров, но плохо масштабируется на большие выборки.

\textbf{LOF (Local Outlier Factor)} вычисляет относительную плотность точки
по сравнению с её \(k\)-ближайшими соседями. Точка с локальной плотностью,
существенно ниже плотности соседей, считается выбросом.

\textbf{Isolation Forest}~\cite{liu2008isolation} использует случайное
рекурсивное разбиение пространства; чем меньше глубина, на которой точка
изолируется, тем выше её аномальность. Для узла-листа с глубиной \(h(x)\)
скоринг определяется как
\begin{equation}
s(x, n) = 2^{-\frac{\mathbb{E}[h(x)]}{c(n)}},
\qquad
c(n) \approx 2\bigl(\ln(n-1) + 0.5772\bigr) - \frac{2(n-1)}{n},
\label{eq:iforest}
\end{equation}
где \(c(n)\) --- средняя длина пути в случайном бинарном дереве из \(n\)
точек. Метод отличается линейной по объёму данных сложностью и хорошо
переносится на признаковые пространства NIDS.

\subsection{One-class классификация}

\textbf{One-Class SVM}~\cite{scholkopf2001ocsvm} ищет в спрямляющем
пространстве минимальную гиперсферу, содержащую долю \(1-\nu\) обучающих
наблюдений, что эквивалентно построению разделяющей гиперплоскости с
максимальным отступом от начала координат:
\begin{equation}
\min_{w, \rho, \xi} \tfrac{1}{2}\|w\|^{2} + \tfrac{1}{\nu n}\sum_{i=1}^{n}\xi_i - \rho,
\quad \text{s.t.}\quad \langle w, \phi(x_i)\rangle \geq \rho - \xi_i,\ \xi_i \geq 0,
\label{eq:ocsvm}
\end{equation}
где \(\phi\) --- спрямляющее преобразование, индуцированное ядром.
Аномалия идентифицируется как отрицательное значение функции
\(f(x) = \langle w, \phi(x)\rangle - \rho\).

Параметр \(\nu \in (0, 1]\) задаёт верхнюю оценку доли выбросов в обучающей
выборке. Метод чувствителен к выбору ядра (радиальное гауссовское ядро
наиболее распространено) и требует тщательного масштабирования признаков.

\subsection{Нейросетевые методы детекции аномалий}

Нейросетевые модели в задачах NIDS реализуются в одном из двух режимов:
дискриминативном (бинарная или мультиклассовая классификация) и
реконструкционно-вероятностном (автоэнкодеры, вариационные автоэнкодеры).
Принципиальное отличие от классических методов состоит в способности
автоматически извлекать иерархические признаки из исходных flow-векторов и
последовательностей соединений~\cite{lecun2015dl, vinayakumar2019dl,
mirsky2018kitsune}. Подробный анализ архитектур и обоснование выбора --- в
\secref{sec:dl-architectures}.

\section{Глубинное обучение для анализа сетевого трафика}
\label{sec:dl-architectures}

\subsection{Полносвязные нейронные сети (MLP)}

Многослойный перцептрон выступает базовой моделью DL и описывается
рекурсией
\begin{equation}
h^{(l)} = \phi\!\left(W^{(l)} h^{(l-1)} + b^{(l)}\right),
\qquad l = 1, \dots, L,
\label{eq:mlp}
\end{equation}
где \(\phi\) --- нелинейность (ReLU, LeakyReLU, GELU), \(W^{(l)}\), \(b^{(l)}\)
--- обучаемые параметры. Для бинарной классификации выход формируется как
\(\hat{p}(y=1\mid x) = \sigma(W^{(L+1)} h^{(L)} + b^{(L+1)})\). Обучение
сводится к минимизации бинарной перекрёстной энтропии
\begin{equation}
\mathcal{L}_{\text{BCE}} =
-\frac{1}{n}\sum_{i=1}^{n}\Bigl[y_i \log \hat{p}_i + (1-y_i)\log(1-\hat{p}_i)\Bigr].
\label{eq:bce}
\end{equation}
В контексте NIDS MLP служит нижней границей DL-моделей: он не учитывает ни
локальной, ни временной структуры, но позволяет проверить достаточность
признакового представления.

\subsection{Свёрточные сети для последовательностей (1D-CNN)}

Одномерные свёрточные слои применяются для извлечения локальных паттернов из
последовательностей flow-признаков в скользящем окне. Операция
свёртки имеет вид
\begin{equation}
(h * w)_t = \sum_{k=0}^{K-1} w_k\, h_{t-k},
\label{eq:conv1d}
\end{equation}
где \(K\) --- размер ядра. Стек свёрточных слоёв с пуллингом формирует
иерархию рецептивных полей, позволяя обнаруживать локальные ``мотивы''
поведения, характерные для конкретных классов атак (например, пакеты SYN-флуда
в окне порядка десятков соединений). По данным
работ~\cite{vinayakumar2019dl, ferrag2020dlcyber}, 1D-CNN превосходит MLP по
полноте детекции атак типа DoS и Reconnaissance.

\subsection{Рекуррентные сети: RNN, LSTM, GRU, BiLSTM}

Рекуррентная сеть моделирует временные зависимости через скрытое состояние
\(h_t\):
\begin{equation}
h_t = \phi(W_{xh} x_t + W_{hh} h_{t-1} + b_h).
\label{eq:rnn}
\end{equation}
Простая RNN страдает от проблемы исчезающего/взрывающегося градиента, что
делает её малопригодной для длинных последовательностей.

\textbf{LSTM}~\cite{hochreiter1997lstm} вводит ячейку состояния
\(c_t\) и три гейта (входной \(i_t\), забывания \(f_t\), выходной \(o_t\)):
\begin{align}
i_t &= \sigma(W_i [h_{t-1}, x_t] + b_i),\\
f_t &= \sigma(W_f [h_{t-1}, x_t] + b_f),\\
o_t &= \sigma(W_o [h_{t-1}, x_t] + b_o),\\
\tilde{c}_t &= \tanh(W_c [h_{t-1}, x_t] + b_c),\\
c_t &= f_t \odot c_{t-1} + i_t \odot \tilde{c}_t,\\
h_t &= o_t \odot \tanh(c_t).
\label{eq:lstm}
\end{align}
Гейтовая структура обеспечивает устойчивое распространение градиента и
способность модели хранить долгосрочный контекст --- свойство, востребованное
при обнаружении медленных атак (\eng{slow brute-force}, \eng{low-and-slow
exfiltration}).

\textbf{GRU}~\cite{cho2014gru} представляет собой упрощённый вариант с двумя
гейтами (\(r_t\), \(z_t\)) и без отдельной ячейки состояния. По числу
параметров GRU экономичнее LSTM при сопоставимом качестве на коротких
последовательностях.

\textbf{BiLSTM} обучает два LSTM в прямом и обратном направлениях, что
позволяет учитывать контекст ``из будущего'' при обработке элемента в
середине окна. Применимо к режимам пост-фактум анализа; для
real-time-инференса требует обработки полного окна.

\subsection{Гибридные модели CNN--LSTM}

Комбинация одномерной свёртки с рекуррентным слоем сочетает преимущества обеих
архитектур: CNN извлекает локальные паттерны (\eng{spatial features}) в
скользящем окне flow-признаков, LSTM моделирует временные зависимости между
получаемыми CNN представлениями~\cite{vinayakumar2019dl, kim2016lstm}. Схема
обработки имеет вид
\[
x_t \;\xrightarrow{\,\text{Conv1D}\,}\; z_t \;\xrightarrow{\,\text{LSTM}\,}\; h_t
\;\xrightarrow{\,\text{FC}\,}\; \hat{p}_t.
\]
Такая комбинация согласуется с природой сетевого трафика, в котором аномалия
может проявляться одновременно как ``локальный шаблон'' (последовательность
соседних соединений) и ``временной тренд'' (продолжительная активность). По
этой причине CNN--LSTM выбран в настоящей работе в качестве основной
\eng{proposed}-архитектуры (детальное обоснование --- в главе~2).

\subsection{TCN и Transformer-кодировщики}

\textbf{TCN (Temporal Convolutional Network)}~\cite{bai2018tcn} использует
причинные дилатированные свёртки и полностью обходится без рекуррентных
связей; рецептивное поле растёт экспоненциально по числу слоёв,
\(R = 1 + (K-1)\sum_{l=0}^{L-1} 2^l\). Преимущества: высокая параллелизуемость,
устойчивость градиента. Недостаток: фиксированное рецептивное поле, что в
NIDS может приводить к занижению длинных зависимостей.

\textbf{Transformer-encoder}~\cite{vaswani2017attention} основан на
механизме само-внимания
\begin{equation}
\operatorname{Attention}(Q, K, V) =
\operatorname{softmax}\!\left(\frac{QK^{\top}}{\sqrt{d_k}}\right) V,
\label{eq:attention}
\end{equation}
где \(Q\), \(K\), \(V\) --- проекции входной последовательности.
Само-внимание обеспечивает прямой доступ ко всему окну и эффективное
моделирование длинных зависимостей. Применение к NIDS требует учёта позиции
(\eng{positional encoding}) и существенно зависит от объёма обучающих данных:
по данным обзоров~\cite{ferrag2020dlcyber, ahmad2021nids}, для UNSW-NB15
TT-модели показывают результат, сопоставимый с CNN--LSTM, но с большей
дисперсией по seed.

\subsection{Автоэнкодеры и вариационные автоэнкодеры}

Автоэнкодер обучается восстанавливать вход \(x\) через узкое горлышко
\(z = E(x)\):
\begin{equation}
\mathcal{L}_{\text{AE}} = \frac{1}{n}\sum_{i=1}^{n} \|x_i - D(E(x_i))\|^{2}.
\label{eq:ae}
\end{equation}
Аномалия идентифицируется как наблюдение с ошибкой реконструкции, превышающей
порог. На таком принципе построен Kitsune~\cite{mirsky2018kitsune} ---
эталонная online-модель ансамбля автоэнкодеров для NIDS.

Вариационный автоэнкодер (VAE)~\cite{kingma2014vae} формализует кодирование
как вероятностный процесс
\begin{equation}
\mathcal{L}_{\text{VAE}} =
\mathbb{E}_{q(z|x)}[\log p(x|z)] - \KL\!\left(q(z|x)\,\|\,p(z)\right),
\label{eq:vae}
\end{equation}
где \(p(z)\) --- априорное распределение скрытого пространства,
\(q(z|x)\) --- вариационный апостериор. VAE позволяет получать вероятностный
скоринг и оценивать неопределённость предсказаний; в NIDS-применениях
демонстрирует устойчивость к шумам в обучающей выборке~\cite{lopez2017cvae}.

\subsection{Сводное теоретическое сравнение архитектур}

Сводное сопоставление, основанное на теоретических свойствах и
экспериментальных данных, опубликованных в обзорных
работах~\cite{khraisat2019survey, ahmad2021nids, liu2019survey,
ferrag2020dlcyber, vinayakumar2019dl}, представлено в \tabref{tab:dl-arch-compare}.
Численные диапазоны метрик не приводятся, поскольку зависят от датасета,
схемы оценки и не должны вырываться из контекста; раздел носит качественный
характер. Количественное сравнение выполняется в главе~2 на основании
повторяемых экспериментов.

\begin{table}[H]
\centering
\caption{Качественное сравнение нейросетевых архитектур для NIDS}
\label{tab:dl-arch-compare}
\renewcommand{\arraystretch}{1.2}
\begin{tabularx}{\linewidth}{@{}lXXXX@{}}
\toprule
\textbf{Модель} & \textbf{Локальные паттерны} & \textbf{Временные зависимости} &
\textbf{Сложность обучения} & \textbf{Inference latency} \\
\midrule
MLP        & нет          & нет        & низкая  & низкая \\
1D-CNN     & да           & ограничено & средняя & низкая \\
LSTM       & нет          & да         & средняя & средняя \\
GRU        & нет          & да         & средняя & низкая \\
BiLSTM     & нет          & двунаправ. & высокая & средняя \\
CNN--LSTM  & да           & да         & высокая & средняя \\
TCN        & да           & фиксиров.  & средняя & низкая \\
Transformer& через attention & да      & высокая & средняя \\
AE/VAE     & косвенно     & нет        & средняя & низкая \\
\bottomrule
\end{tabularx}
\end{table}

Из таблицы следует, что CNN--LSTM покрывает оба измерения структуры flow-данных
(локальное и временное) при разумной вычислительной стоимости. Дополнительные
архитектуры (BiLSTM, Transformer-encoder, TCN) рассматриваются как
конкурентные, и их прямое сопоставление приводится в главе~2.

\section{Метрики качества и протоколы оценки}
\label{sec:metrics}

\subsection{Базовые классификационные метрики}

Для бинарной задачи введём матрицу ошибок (\eng{confusion matrix}) с величинами
\(TP\), \(TN\), \(FP\), \(FN\). Тогда
\begin{align}
\text{Precision} &= \frac{TP}{TP + FP}, \label{eq:precision}\\
\text{Recall (TPR)} &= \frac{TP}{TP + FN}, \label{eq:recall}\\
\text{F1} &= \frac{2 \cdot \text{Precision} \cdot \text{Recall}}
                 {\text{Precision} + \text{Recall}},\label{eq:f1}\\
\text{FPR} &= \frac{FP}{FP + TN}.\label{eq:fpr}
\end{align}

\textbf{Precision} измеряет долю истинных атак среди вызвавших алерт; имеет
прямую интерпретацию в контексте SOC --- доля ``полезных'' алертов.
\textbf{Recall} (он же TPR) измеряет долю обнаруженных атак среди всех
реальных. \textbf{F1} --- гармоническое среднее, удобное при дисбалансе
классов. \textbf{FPR} --- долю неверно сработавших правил среди всех
нормальных наблюдений.

В условиях NIDS дополнительно используется метрика
\(\text{FP-per-time}\) --- ожидаемое число ложных срабатываний в единицу
времени, более показательное при больших объёмах нормального трафика.

Коэффициент \textbf{MCC} (Matthews correlation coefficient) учитывает все
четыре элемента матрицы ошибок и рекомендуется при сильном дисбалансе:
\begin{equation}
\text{MCC} =
\frac{TP \cdot TN - FP \cdot FN}
{\sqrt{(TP+FP)(TP+FN)(TN+FP)(TN+FN)}}.
\label{eq:mcc}
\end{equation}

\subsection{ROC-AUC, PR-AUC и кривые качества}

Поскольку детектор формирует непрерывный скоринг \(s(x)\), его качество
характеризуется не только в одной рабочей точке, но и интегрально.

\textbf{ROC-AUC} --- площадь под кривой, отображающей TPR от FPR при
изменении порога \(\tau\). Для несбалансированных задач даёт оптимистическую
оценку~\cite{khraisat2019survey, ahmad2021nids}.

\textbf{PR-AUC} --- площадь под кривой Precision--Recall. Она более
чувствительна к редкому классу и поэтому считается предпочтительной для
NIDS-сценариев с малой долей атак~\cite{liu2019survey}. В работе
PR-AUC используется как ключевая интегральная метрика качества модели.

\subsection{Эксплуатационные метрики}

К эксплуатационным метрикам относятся:
\begin{itemize}
    \item \textbf{Time-to-detect (TTD)} --- время с момента появления атаки до
    срабатывания алерта, выраженное в единицах окон или в физическом времени;
    \item \textbf{Inference latency} --- время выработки скоринга на одно
    окно;
    \item \textbf{Throughput} --- число обработанных событий в секунду
    (EPS);
    \item \textbf{Resource footprint} --- потребление CPU/GPU/RAM.
\end{itemize}
Эти величины определяют применимость модели в SOC и подбираются под целевую
эксплуатационную нагрузку.

\subsection{Калибровка и приоритизация}

Большинство нейросетевых моделей дают плохо калиброванные вероятности
\(\hat{p}\)~\cite{guo2017calibration}: уверенность модели систематически
смещена. Для использования \(\hat{p}\) в формировании приоритетов алертов
применяется температурное масштабирование
\begin{equation}
\hat{p}_{T}(y=1\mid x) =
\sigma\!\left(\frac{z(x)}{T}\right),
\label{eq:temp-scaling}
\end{equation}
где \(z(x)\) --- логит, \(T\) --- параметр, подбираемый на валидации.
Для оценки качества калибровки используется \textbf{ECE (Expected Calibration
Error)}~\cite{guo2017calibration}.

\subsection{Метрики дрейфа распределений}

Для контроля стабильности входного распределения используются:

\begin{itemize}
    \item \textbf{PSI (Population Stability Index)}: при разбиении значений
    признака на бины \(b\)
    \begin{equation}
    \PSI = \sum_{b}\,(p_b - q_b)\ln\frac{p_b}{q_b},
    \label{eq:psi}
    \end{equation}
    где \(p_b\), \(q_b\) --- доли наблюдений в бине \(b\) для эталонного и
    нового окон;
    \item \textbf{KL-расхождение}
    \(\KL(P\|Q) = \sum_b p_b \ln \tfrac{p_b}{q_b}\);
    \item \textbf{MMD (Maximum Mean Discrepancy)} --- мера различия между
    распределениями в воспроизводящем гильбертовом пространстве,
    \begin{equation}
    \MMD^{2}(P, Q) =
    \mathbb{E}_{x,x'\sim P}[k(x,x')] - 2\mathbb{E}_{x\sim P, y\sim Q}[k(x,y)]
    + \mathbb{E}_{y,y'\sim Q}[k(y,y')].
    \label{eq:mmd}
    \end{equation}
\end{itemize}

В рамках экспериментов главы~4 указанные метрики применяются для оценки
эффекта drift-инжектора и стабильности модели на потоковом сценарии.

\section{Ограничения существующих решений и постановка задачи исследования}

Системный анализ публикаций последних лет выявил следующие нерешённые
проблемы, релевантные для настоящей работы.

\begin{enumerate}
    \item Большинство экспериментов в литературе опирается на классические
    датасеты (KDD99, NSL-KDD), не отражающие современные классы атак, что
    приводит к завышенным оценкам качества~\cite{tavallaee2009nslkdd,
    moustafa2016eval, sommer2010outside, gharib2016evalframework}.
    \item Сравнение моделей часто выполняется без явного контроля дрейфа
    распределений и без оценки устойчивости при смещении тестовой выборки
    относительно обучающей~\cite{pendlebury2019tesseract,
    andresini2021insomnia}.
    \item Отсутствует общедоступный воспроизводимый стенд, в рамках которого
    модель детекции и активный (генеративный) источник атак работают в режиме
    реального времени и допускают параметризованное переключение режимов
    ``норма/атака/дрейф''~\cite{ring2019flowgan, apruzzese2020adversarial}.
    \item Большая часть исследований сосредоточена на оптимизации бинарной
    классификации и не учитывает требований SOC к приоритизации алертов и
    стоимости ложных срабатываний на единицу времени~\cite{sommer2010outside}.
\end{enumerate}

Перечисленные пробелы определяют целевые задачи диссертационной работы:
обоснованный выбор архитектуры на основе теоретического и экспериментального
сравнения; построение методики обнаружения, явно учитывающей дрейф
распределений; реализация ИИ-злоумышленника в форме воспроизводимого
шаблонного агента с конечно-автоматной политикой режимов.

\section{Промежуточные выводы по разделу}

\begin{enumerate}
    \item Корпоративная сеть характеризуется неоднородностью трафика,
    дисбалансом легитимного и атакующего поведения и динамичностью
    распределений. Это диктует необходимость использовать flow-уровень
    представления данных и модели, способные учитывать совместную и временную
    структуру признаков.
    \item Сигнатурный, аномальный и гибридный подходы NIDS дополняют, а не
    заменяют друг друга. Аномальный подход на основе нейросетей выбран в
    работе по причине способности обнаруживать ранее неизвестные сценарии
    при принципиальной возможности интеграции с сигнатурной ветвью на
    уровне SOC.
    \item Среди методов обнаружения аномалий нейросетевые модели обладают
    наибольшим потенциалом за счёт автоматического выделения иерархических
    признаков. Среди архитектур наиболее адекватной для flow-последовательностей
    является CNN--LSTM, что согласуется с теоретическим анализом и обзорной
    литературой и принимается в качестве основной (\eng{proposed}) модели.
    \item Корректная оценка качества требует одновременного использования
    интегральных (PR-AUC, MCC), пороговых (Precision, Recall, F1) и
    эксплуатационных (FP-per-time, time-to-detect, latency) метрик; для
    оценки устойчивости введены метрики дрейфа PSI, KL и MMD.
    \item Систематизация ограничений существующих решений подтверждает
    актуальность задач, заявленных в настоящей работе: построение методики,
    объединяющей нейросетевую детекцию с воспроизводимым активным
    тестированием и явной оценкой устойчивости к дрейфу.
\end{enumerate}
